\documentclass[a4paper, 10.5pt]{article}

% ── PACKAGES ─────────────────────────────────────────────
\usepackage[utf8]{inputenc}
\usepackage[T1]{fontenc}
\usepackage{lmodern}
\usepackage[left=1.7cm, right=1.7cm, top=1.4cm, bottom=1.4cm]{geometry}
\usepackage[hidelinks]{hyperref}
\usepackage{fontawesome5}
\usepackage{enumitem}
\usepackage{xcolor}
\usepackage{titlesec}
\usepackage{setspace}

% ── COLORS ───────────────────────────────────────────────
\definecolor{heading}{RGB}{30, 30, 30}
\definecolor{subtext}{RGB}{80, 80, 80}
\definecolor{accent}{RGB}{0, 70, 140}
\definecolor{rulecolor}{RGB}{60, 60, 60}

% ── PAGE SETUP ───────────────────────────────────────────
\pagestyle{empty}
\raggedright
\setlength{\parindent}{0pt}
\setlength{\parskip}{0pt}
\setstretch{1.08}

% ── SECTION FORMAT ───────────────────────────────────────
\titleformat{\section}
  {\large\bfseries\color{heading}\scshape}
  {}{0em}{}
  [\vspace{-6pt}{\color{rulecolor}\rule{\textwidth}{0.6pt}}\vspace{-2pt}]
\titlespacing*{\section}{0pt}{8pt}{4pt}

% ── LIST FORMAT ──────────────────────────────────────────
\setlist[itemize]{
  noitemsep,
  topsep=2pt,
  parsep=1pt,
  leftmargin=1.4em,
  label={\small\textbullet}
}

% ── CUSTOM COMMANDS ──────────────────────────────────────
\newcommand{\entry}[3]{%
  \vspace{4pt}%
  \textbf{#1} \hfill \textcolor{subtext}{\small #3}\\[-2pt]
  \textit{\textcolor{subtext}{#2}}%
  \vspace{2pt}%
}

\newcommand{\entrysingle}[2]{%
  \vspace{4pt}%
  \textbf{#1} \hfill \textcolor{subtext}{\small #2}%
  \vspace{2pt}%
}

% ══════════════════════════════════════════════════════════
\begin{document}

% ── HEADER ───────────────────────────────────────────────
\begin{center}
  {\LARGE\bfseries\color{heading} Ashish Joshi}\\[7pt]
  {\small
    \faEnvelope\ \href{mailto:connectwithashish@gmail.com}{connectwithashish@gmail.com}
    \enspace\textbar\enspace
    \faPhone\ +91 9305292142
    \enspace\textbar\enspace
    \faLinkedin\ \href{https://www.linkedin.com/in/ashish-joshi-6204442b7/}{LinkedIn}
    \enspace\textbar\enspace
    \faGithub\ \href{https://github.com/31ASHISH}{GitHub}
    \enspace\textbar\enspace
    \faCode\ \href{https://leetcode.com/u/Ashish_Joshi31/}{LeetCode}
  }
\end{center}

\vspace{-2pt}

% ── EDUCATION ────────────────────────────────────────────
\section{Education}

\textbf{Indian Institute of Technology, Palakkad} \hfill \textcolor{subtext}{\small Aug 2023 -- Present}\\
B.Tech in Civil Engineering \hfill CGPA: \textbf{8.60}\\[2pt]
{\textbf{AI \& Data Science Focus:} Machine Learning, Deep Learning, NLP, Agentic AI, Large Language Models}\\
{\textbf{Technical Exposure:} Statistical modeling, spatiotemporal datasets, data-driven analysis}

% ── PROFILE ──────────────────────────────────────────────
\section{Profile}

{\small
AI-focused undergraduate with hands-on experience building \textbf{LLM-powered agentic systems}, RAG pipelines, and deep learning models across NLP, climate science, signal processing, and financial analytics.
Driven to contribute to \textbf{high-impact, publishable research} in Agentic AI, Generative AI, and Large Language Models.
}

% ── TECHNICAL SKILLS ─────────────────────────────────────
\section{Technical Skills}

\textbf{AI \& Machine Learning:} Generative AI, Large Language Models (LLMs), Agentic AI, Retrieval-Augmented Generation (RAG), NLP, Deep Learning, Time-Series \& Spatiotemporal Modeling\\[3pt]
\textbf{Frameworks \& Libraries:} PyTorch, Hugging Face Transformers, TensorFlow, Scikit-learn\\[3pt]
\textbf{Programming Languages:} Python, C++\\[3pt]
\textbf{Data \& Analysis:} Feature Engineering, Exploratory Data Analysis (EDA), Statistical Evaluation, Dataset Curation\\[3pt]
\textbf{Developer Tools:} Git, GitHub, Jupyter Notebook, Google Colab, Linux, Kaggle, LaTeX

% ── EXPERIENCE ───────────────────────────────────────────
\section{Experience}

\entry{IPTIF Research Intern}{IIT Palakkad}{Feb 2026 -- Present}
\begin{itemize}
  \item Developing a hybrid drought prediction framework integrating \textbf{Wavelet Transform with 4 ML/DL models} (CNN, XGBoost, ConvLSTM, WCNN), analyzing \textbf{5,000+ lat/lon grid points} across India over \textbf{1950--2023} monthly climate records.
  \item Spearheading a \textbf{patent-oriented Wavelet--ML system} combining multi-scale signal decomposition with learning-based prediction for national-scale drought index forecasting.
\end{itemize}

\entry{EAYL Innovation Lab Intern}{IIT Palakkad}{Feb 2026 -- Present}
\begin{itemize}
  \item Architecting the database schema and backend workflows for the \textbf{IITPKD campus application}, designing modular API endpoints with scalable data integration.
  \item Automating manual administrative workflows using \textbf{n8n}, incorporating global search capabilities and event-driven application features.
\end{itemize}

\entry{Data Analytics Intern}{MSME Sampark, India}{May 2025 -- Jul 2025}
\begin{itemize}
  \item Developed Python-based analytical pipelines to process structured datasets from multiple early-stage startups, automating data cleaning and feature extraction.
  \item Deployed statistical models and EDA dashboards to surface actionable insights, informing strategic recommendations for partner startups.
\end{itemize}

\entry{Product Developer Intern (Data \& AI Focus)}{Plant Khareedo, India}{May 2024 -- Jul 2024}
\begin{itemize}
  \item Analyzed customer transaction data to quantify demand patterns, purchase behavior, and user engagement trends across product categories.
  \item Executed A/B experiments evaluating growth strategies and partnered with product teams to implement data-informed feature prioritization.
\end{itemize}

\newpage

% ── PROJECTS ─────────────────────────────────────────────
\section{Projects}

\entrysingle{Real-Time Financial Intelligence System}{\normalfont\small\textcolor{subtext}{Python, LLMs, RAG, LSTM, Agentic AI}}
\begin{itemize}
  \item Architected an \textbf{agentic AI system} fusing financial time-series, SEC filings, and market news with RAG-based retrieval, evaluated on \textbf{FinBen}, \textbf{FLUE}, and \textbf{FinanceBench} benchmarks.
  \item Implemented RAG pipelines enabling LLMs to perform contextual reasoning over financial knowledge bases, benchmarked on \textbf{FPB} sentiment and \textbf{FiQA} opinion mining tasks.
  \item Integrated LSTM-based time-series forecasting with LLM-driven agentic reasoning for automated portfolio risk assessment and decision-support.
\end{itemize}

\entrysingle{SPI Prediction using ML/DL Models for Drought Forecasting}{\normalfont\small\textcolor{subtext}{Python, ML/DL, Spatiotemporal Modeling}}
\begin{itemize}
  \item Investigated drought predictability across \textbf{70 lat/lon points} in India using monthly precipitation data from \textbf{1950--2021} and the Standardized Precipitation Index (SPI).
  \item Analyzed the influence of climate oscillation indices (ENSO, PDO, IOD) through correlation analysis and feature engineering to identify key drought predictors.
  \item Trained and benchmarked CNN, RNN, ConvLSTM, XGBoost, and Wavelet-CNN models to capture nonlinear spatiotemporal dependencies in SPI prediction.
  \item Assessed model generalization using comparative error metrics and temporal cross-validation across different Indian climate zones.
\end{itemize}

\entrysingle{Large-Scale Sentiment Analysis on Amazon Kindle Reviews}{\normalfont\small\textcolor{subtext}{NLP, Representation Learning}}
\begin{itemize}
  \item Processed \textbf{980K+ customer reviews} through NLP preprocessing pipelines including tokenization, stopword removal, lemmatization, and vectorization.
  \item Benchmarked classical ML models (Logistic Regression, SVM, Random Forest) against LSTM and BiLSTM architectures for sentiment classification.
  \item Compared TF-IDF and Word2Vec embedding strategies to evaluate dense vs.\ sparse representation impact on deep learning sequence models.
\end{itemize}

\entrysingle{Spectrum-SLM: Small Language Model for Cognitive Radio Spectrum Sensing}{\normalfont\small\textcolor{subtext}{PyTorch, Transformers, SDR, Signal Processing}}
\begin{itemize}
  \item Designed a \textbf{GPT-inspired Small Language Model ($\sim$1M params)} for real-time spectrum sensing using ADALM-Pluto SDR, treating \textbf{176-bin} PSD vectors as token sequences.
  \item Curated a \textbf{152K+ sample} multi-modulation dataset (BPSK, QPSK, 8PSK, 16QAM) from real RF measurements at 2.4\,GHz with GNU Radio.
  \item Engineered a patch-based spectrum tokenizer with a 3-phase pipeline: self-supervised Masked Spectrum Modeling, multi-task fine-tuning, and autoregressive prediction.
  \item Benchmarked against ensemble baselines (RandomForest + XGBoost + AdaBoost), targeting improved low-SNR detection ($<$8\,dB) with multi-task inference.
\end{itemize}

% ── ACHIEVEMENTS ─────────────────────────────────────────
\section{Achievements}

\begin{itemize}
  \item Awarded \textbf{Honda Hope Scholarship 2024} --- recognized for academic excellence and innovation.
  \item Secured \textbf{5th Place at Engineers' Conclave (Inter-IIT)} --- evaluated for technical depth among teams from 23 IITs.
  \item Achieved \textbf{2nd Place, COGNIFAI Competition} (NLP \& ML).
  \item Ranked \textbf{Top 10, Agentic AI Pathway --- Inter IIT Tech Meet 8.0}.
  \item Won \textbf{1st Place, FinTech Competition}.
  \item Solved \textbf{160+ problems on LeetCode}.
\end{itemize}

% ── INTERESTS ────────────────────────────────────────────
\section{Areas of Interest}

{\small Agentic AI \enspace$\cdot$\enspace Generative AI \enspace$\cdot$\enspace Large Language Models \enspace$\cdot$\enspace Retrieval-Augmented Generation \enspace$\cdot$\enspace AI Agents \enspace$\cdot$\enspace Time-Series and Spatiotemporal Modeling}

\end{document}
